% Part: computability
% Chapter: computability-theory
% Section: halting-problem

\documentclass[../../../include/open-logic-section]{subfiles}

\begin{document}

\olfileid{cmp}{thy}{hlt}
\olsection{અટકવાની સમસ્યા}

રચના મુજબ, સર્વવ્યાપી આંશિક સંગણનીય વિધેય~$\fn{Un}(e,x)$ ત્યારે અને
માત્ર ત્યારે વ્યાખ્યાયિત છે જ્યારે~$e$ વડે કોડ કરેલા વિધેયની
સંગણના ઇનપુટ~$x$ માટે મૂલ્ય આપે. આવું થાય છે કે નહીં તે નક્કી
કરી શકીએ કે કેમ એવો પ્રશ્ન સ્વાભાવિક છે. હકીકતમાં તે શક્ય નથી.
સંગણનાના ટ્યુરિંગ મશીન નિદર્શમાં તેનો અર્થ એ થાય કે આપેલું
ટ્યુરિંગ મશીન આપેલા ઇનપુટ પર થંભશે કે નહીં તે સંગણકીય રીતે
અનિર્ણેય છે. આથી આગળનું પ્રમેય ``અટકવાની સમસ્યાની અનિર્ણેયતા''
તરીકે ઓળખાય છે. નીચે બે સાબિતીઓ આપીશું. પહેલી અગાઉની ચર્ચાને
આગળ ધપાવે છે અને બીજી વધુ સીધી છે.

\begin{thm}
\ollabel{thm:halting-problem}
માનો કે
\[
h(e, x)  =
\begin{cases}
1 & \text{જો\/ $\fn{Un}(e, x)$ વ્યાખ્યાયિત હોય} \\
0 & \text{અન્યથા.}
\end{cases}
\]
તો $h$ સંગણનીય નથી.
\end{thm}

\begin{proof}
માનો કે $h$ સંગણનીય છે. તો તેનાથી સર્વવ્યાપી સંપૂર્ણ સંગણનીય
વિધેય વ્યાખ્યાયિત કરી શકાય એવું બતાવીશું. વ્યાખ્યાયિત કરો:
\[
\fn{Un'}(e,x) =
\begin{cases}
\fn{Un}(e,x) & \text{જો $h(e,x) = 1$} \\
0 & \text{અન્યથા.}
\end{cases}
\]
હવે $\fn{Un'}(e, x)$ સંપૂર્ણ વિધેય છે અને $h$ સંગણનીય હોય તો
તે પણ સંગણનીય છે. ઉદાહરણ તરીકે, નીચે પ્રમાણે આદિમ પુનરાવર્તનથી
$g$ વ્યાખ્યાયિત કરી શકાય:
\begin{align*}
g(0, e, x) & \simeq 0 \\
g(y+1, e, x) & \simeq \fn{Un}(e,x);
\end{align*}
ત્યારે
\[
\fn{Un'}(e,x) \simeq g(h(e,x),e,x).
\]
જ્યાં $\fn{Un}(e,x)$ વ્યાખ્યાયિત છે ત્યાં $\fn{Un'}(e,x)$ તેની
સાથે મેળ ખાય છે. આથી જે આંશિક સંગણનીય વિધેયો હકીકતમાં સંપૂર્ણ
છે, તેમના માટે $\fn{Un'}$ સર્વવ્યાપી છે. પરંતુ આ
\olref[nou]{thm:no-univ} સાથે વિરોધાભાસ છે.
\end{proof}

\begin{proof}
માનો કે $h(e,x)$ સંગણનીય છે. વિધેય $g$ આ રીતે વ્યાખ્યાયિત કરો:
\[
g(x) =
\begin{cases}
  0                & \text{જો $h(x,x) = 0$} \\
  \fundefined & \text{અન્યથા.}
\end{cases}
\]
$g$ આંશિક સંગણનીય છે. દાખલા તરીકે, તેને
$\umin{y}{h(x,x) = 0}$ તરીકે વ્યાખ્યાયિત કરી શકાય. તેથી કોઈ~$e$
માટે દરેક~$x$ પર $g(x) \simeq \fn{Un}(e,x)$ છે. શું $g$, $e$ પર
વ્યાખ્યાયિત છે? જો હોય, તો $g$ની વ્યાખ્યા મુજબ $h(e,e) = 0$ છે
($h$ વ્યાખ્યાયિત હોય ત્યારે જ~$0$ મૂલ્ય લઈ શકે છે).
$h$ની વ્યાખ્યા મુજબ તેથી $\fn{Un}(e, e)$ અનિર્ધારિત છે.
દરેક~$x$ માટે $g(x) \simeq \fn{Un}(e, x)$ હોવાની આપણી ધારણા
મુજબ $g(e)$ અનિર્ધારિત થાય, જે વિરોધાભાસ છે. બીજી તરફ, જો
$g(e)$ અનિર્ધારિત હોય, તો $h(e,e) \neq 0$, એટલે
$h(e,e) = 1$. તેથી $\fn{Un}(e, e)$ વ્યાખ્યાયિત છે. પરંતુ
$g(x) \simeq \fn{Un}(e, x)$ હોવાથી $g(e)$ પણ વ્યાખ્યાયિત
હોવું જોઈએ. આ પણ વિરોધાભાસ છે.
\end{proof}

\begin{tagblock}{TMs}
\begin{explain}
આ દલીલ ટ્યુરિંગ મશીનોની ભાષામાં પણ કહી શકાય. માનો કે એવું
ટ્યુરિંગ મશીન~$H$ છે જે ટ્યુરિંગ મશીન~$E$નું વર્ણન અને ઇનપુટ~$x$
લઈને, $E$ ઇનપુટ~$x$ પર થંભે છે કે નહીં તે નક્કી કરે છે.
તો એક ઇનપુટ~$x$ લેતું બીજું ટ્યુરિંગ મશીન~$G$ બનાવી શકાય:
તે $H$ ચલાવીને તપાસે કે સૂચક~$x$ ધરાવતું મશીન $M_x$
ઇનપુટ~$x$ પર થંભે છે કે નહીં, અને પછી વિરુદ્ધ કરે. એટલે કે,
જો $H$ કહે કે $M_x$ ઇનપુટ~$x$ પર થંભે છે, તો $G$ અનંત લૂપમાં
જાય છે; અને જો $H$ કહે કે $M_x$ ઇનપુટ~$x$ પર થંભતું નથી,
તો $G$ થંભી જાય છે. પોતાના સૂચકને ઇનપુટ તરીકે આપીએ ત્યારે
$G$ થંભે છે? ઉપરની દલીલ કહે છે કે તે ત્યારે અને માત્ર ત્યારે
થંભે છે જ્યારે તે થંભતું નથી---વિરોધાભાસ. તેથી એવું ટ્યુરિંગ
મશીન~$H$ છે એવી આપણી ધારણા ખોટી છે.
\end{explain}
\end{tagblock}

\end{document}
